\subsection{Approximation for essential valuations}

As the essential valuations of a Krull domain are discrete and normed, no two of them are equivalent and hence they are independent (Chapter VI, § 7, no. 2). Corollary 2 to the approximation theorem (loc. cit., Theorem 1) may therefore be applied to them: given some \(n_i \in \mathbf{Z}\) and some essential valuations \(v_i\) finite in number and distinct, there exists \(x \in \mathbf{K}\) such that \(v_i(x) = n_i\) for all \(i\) . But here there is a more precise result:

PROPOSITION 9. Let \(v_1, \ldots, v_r\) be distinct essential valuations of a Krull domain A and \(n_1, \ldots, n_r\) rational integers. There exists an element \(x\) of the field of fractions K of A such \(v_i(x) = n_i\) for \(1 \leqslant i \leqslant r\) and \(v(x) \geqslant 0\) for every essential valuation \(v\) of A distinct from \(v_1, \ldots, v_r\) .

Let \(\mathfrak{p}_1, \ldots, \mathfrak{p}_r\) be the divisorial ideals of A corresponding to the valuations \(v_1, \ldots, v_r\) . There exists \(y \in K\) such that \(v_i(y) = n_i\) for \(1 \leqslant i \leqslant r\) (Chapter VI, § 7, no. 2, Corollary 1 to Theorem 1). The essential valuations \(w_1, \ldots, w_s\) of A distinct from the \(v_i\) for which the integer \(w_j(y) = -m_j\) is \(<0\) are finite in number; let \(\mathfrak{q}_1, \ldots, \mathfrak{q}_s\) be the corresponding ideals. There exists no inclusion relation between \(\mathfrak{p}_1, \ldots, \mathfrak{p}_r, \mathfrak{q}_1, \ldots, \mathfrak{q}_s\) since these ideals correspond to extremal divisors and these ideals are prime (Corollary 1 to Proposition 6). Hence the integral ideal \(\mathfrak{a} = \mathfrak{q}_1^{m_1} \ldots \mathfrak{q}_s^{m_s}\) is contained in none of the \(\mathfrak{p}_i\) (Chapter II, § 1, no. 1, Proposition 1) and is therefore not contained in their union (loc. cit., Proposition 2). Therefore there exists \(z \in \mathfrak{a}\) such that \(z \notin \mathfrak{p}_i\) for \(1 \leqslant i \leqslant r\) ; then \(v_1(z) = \cdots = v_r(z) = 0\) and \(w_j(z) \geqslant m_j\) for \(1 \leqslant j \leqslant s\) ; hence the element \(x = yz\) solves the problem.

COROLLARY 1. Let A be a Krull domain, K its field of fractions and a, b and c three divisorial fractional ideals of A such that a ⊂ b. There exists x ∈ K such that a = b ∩ xc.

Let \((v_{\iota})_{\iota \in I}\) be the family of essential valuations of A and let \((m_{\iota})\) (resp. \((n_{\iota})\) , \((p_{\iota})\) ) be the family of rational integers (zero except for a finite number of indices) such that \(\mathfrak{a}\) (resp. \(\mathfrak{b}, \mathfrak{c}\) ) is the set of \(x \in K\) for which \(v_{\iota}(x) \geqslant m_{\iota}\) (resp. \(n_{\iota}, p_{\iota}\) ) for all \(\iota \in I\) (Proposition 5, no. 4). The set J of \(\iota \in I\) such that \(m_{\iota} > n_{\iota}\) is finite. As \(p_{\iota} = m_{\iota} = 0\) except for a finite number of indices, Proposition 9 shows that there exists \(x \in K^{*}\) such that \(v_{\iota}(x^{-1}) + m_{\iota} = p_{\iota}\) for \(\iota \in J\) and

\[
v _ {\iota} (x ^ {- 1}) + m _ {\iota} \geqslant p _ {\iota}
\]

for \(\iota \in \mathbf{I} - \mathbf{J}\) . Then, for all \(\iota \in \mathbf{I}\) , \(m_{\iota} = \sup (n_{\iota}, v_{\iota}(x) + p_{\iota})\) . Whence \(\mathfrak{a} = \mathfrak{b} \cap x\mathfrak{c}\) .

COROLLARY 2. Let A be a Krull domain. For a fractional ideal \(\mathfrak{a}\) of A to be divisorial, it is necessary and sufficient that it be the intersection of two fractional principal ideals.

The sufficiency is obvious (no. 1, Definition 2). The necessity follows from Corollary 1: take b and c to be principal and such that b ⊃ a.

